\begin{figure}
	\centering
	\begin{tikzpicture}[scale=0.9]
		% axes
		\draw[thick,->] (0,0) -- (8,0) node[right] {$\bx$};
		\draw[thick,->] (0,0) -- (0,5) node[left] {$y$};
		
		% training data points
		\fill[darkblue] (1,1) circle (3pt);
		\fill[darkblue] (2.5,2.5) circle (3pt);
		\fill[darkblue] (3.1,1.75) circle (3pt);
		\fill[darkblue] (6,3.1) circle (3pt);
		\fill[darkblue] (5,2) circle (3pt);
		\fill[darkblue] (7,4.5) circle (3pt);
		
		% query point
		\node at (4,-0.25) {$\bx_0$};
		\draw[dashed] (4,0) -- (4,5);
		\draw[dashed] (4,2.15) -- (0,2.15) node[left] {$h(\bx_0; \btheta)$};

		% normal curve
		\draw[thick,gray,rounded corners=10pt] (4,0) -- (4,1.2) .. controls (4,1.7) .. (5,2.2) .. controls (4,2.7) .. (4,3.2) -- (4,5)
			node[right,gray] {$n(y \mid \bx_0 ; \btheta, \sigma^2)$};
		\draw[thick,<->] (6,1.9) -- node[right] {$\sigma^2$} (6,2.3);
		
		% regression curve
		\draw[ultra thick,darkblue] (0.5,1) .. controls (3,3) and (6,1) .. (8,4.5) node[right] {$h(\bx; \btheta)$};
	\end{tikzpicture}
	\caption{Visualisierung der probabilistischen Regression. $\bx_0$ ist hierbei ein Testpunkt.}
\end{figure}